\begin{remarks}[7.3.19]
\EGAReaderTarget{ega4.c04.statement.7-3-19}{IV.2.s7.3.19-fr}
\begin{enumerate}
\item[\textup{(i)}] La propiedad $\boldsymbol R$ que sirve para definir
$\boldsymbol P$ satisface la condición $(\mathrm R_{\mathrm{III}})$
en todos los ejemplos enumerados en (7.3.8). Para (i), (ii) y (iii),
esto se deduce de (6.11.2) y del teorema de Cohen (\textbf 0, 19.8.8);
cuando $\boldsymbol P$ es una de las propiedades (iv), (iv$'$), (v),
(v$'$), se deduce de (6.12.7) y (6.12.9), y cuando $\boldsymbol P$
es una de las propiedades (vii), (vii$'$), de (6.12.7) y (6.13.4);
por último, para (vi) y (vi$'$) la afirmación de (7.3.18) es trivial,
pues se cumple para todo preesquema localmente noetheriano.

\oldpage[IV-2]{198}

\item[\textup{(ii)}] Ya hemos señalado (6.4.3) que, cuando
$\boldsymbol P$ es una de las propiedades (ii) o
$(\mathrm{iii}_1)$ de (7.3.8), ignoramos si todo anillo local
noetheriano es un $\boldsymbol P$-anillo; recordemos, sin embargo,
que cuando $A$ es un anillo \emph{cociente de un anillo de
Cohen-Macaulay}, las fibras formales de $A$ son \emph{esquemas de
Cohen-Macaulay} (6.3.8).
\item[\textup{(iii)}] El caso más interesante de la noción de
$\boldsymbol P$-anillo es el que corresponde a la propiedad más
fuerte (iv$'$) de (7.3.8), es decir, a los anillos cuyas
\emph{fibras formales son geométricamente regulares}. Los cuerpos y,
más generalmente, los anillos locales noetherianos \emph{completos}
satisfacen trivialmente esta propiedad.
\item[\textup{(iv)}] Sea $A$ un anillo local noetheriano de
\emph{dimensión 1}; entonces $\operatorname{Spec}(A)$ está formado
por el punto cerrado $a$, correspondiente al ideal maximal
$\mathfrak m$, y por los puntos maximales $b_{\boldsymbol{i}}$
$(1\leqslant i\leqslant r)$, correspondientes a los ideales primos
minimales de $A$. Se tiene $\dim(\widehat A)=1$ y el ideal maximal
de $\widehat A$ es $\mathfrak m\widehat A$; por tanto, la fibra
formal de $A$ en el punto $a$ es $\operatorname{Spec}(k)$, donde
$k=A/\mathfrak m$ es el cuerpo residual de $A$; la fibra formal
en cada uno de los $b_i$ es el espectro de un anillo artiniano
cuyos cuerpos residuales son los cuerpos residuales $L_{ij}$ de
$\operatorname{Spec}(\widehat A)$ en los puntos maximales $b_{\boldsymbol{ij}}$
situados sobre $b_i$. Como un anillo artiniano es un anillo de
Cohen-Macaulay, se ve que $A$ es un $\boldsymbol P$-anillo cuando
$\boldsymbol P$ es la propiedad (ii) de (7.3.8). Además, como un
anillo artiniano reducido es una suma directa de cuerpos, las
propiedades siguientes son equivalentes (6.7.7):
\begin{enumerate}
\item[\emph{a)}] las fibras formales de $A$ son geométricamente
reducidas;
\item[\emph{b)}] las fibras formales de $A$ son geométricamente
normales;
\item[\emph{c)}] las fibras formales de $A$ son geométricamente
regulares;
\end{enumerate}
además, cuando $A$ es \emph{reducido}, también son equivalentes a:
\begin{enumerate}
\item[\emph{d)}] $\widehat A$ es reducido y $L_{\boldsymbol{ij}}$ es una
extensión separable de $K_i$ para todo par $(i,j)$ (4.6.1).
\end{enumerate}
En particular, si $A$ es un \emph{anillo de valoración discreta},
$K$ su cuerpo de fracciones y $\widehat K$ es el completado de $K$
respecto de una valoración correspondiente a $A$ (cuerpo de
fracciones de $\widehat A$), para que las fibras formales de $A$
sean geométricamente regulares es necesario y suficiente que
$\widehat K$ sea una \emph{extensión separable de} $K$. Este será
siempre el caso cuando $K$ sea \emph{de característica} 0.
\end{enumerate}
\end{remarks}
